% !PRINTED 796
\weberpage{796}
\begin{flushleft}
Seite 423, Zeile 3 v.\,u. lies $(-1)^{\frac{(n-1)(n-3)}{4}}=+1$ statt $(-1)^{\frac{n-1}{2}}$. \\
„ 424, Formel (14) „ $P=i^{\frac{n-1}{2}}n^{\frac{n-3}{2}}\sqrt{n}$. \\
„ 436, Zeile 3, 22 v.\,o. lies $\dfrac{\pi}{n}$ statt $\pi$. \\
„ 482, „ 8 v.\,u. ist transitiv und intransitiv zu vertauschen. \\
„ 502, „ 18 v.\,o. lies $\pi_1$ statt $\pi$. \\
„ 504, „ 9 v.\,u. „ $k_{q-1}$ statt $k_{\nu-1}$. \\
„ 514, „ 5 v.\,u. „ $\Theta_1$ statt $\Theta_i$. \\
„ 531, „ 11 v.\,o. „ $y_1^2$ „ $y^2$. \\
„ 540, Formel (10) „ $\gamma_{e-1}$ statt $\gamma_{c-1}$. \\
„ 544, Zeile 15 v.\,o. „ $mB_k=\overset{\varepsilon}{\Sigma}\varepsilon^{-k}(\varepsilon,\alpha)^\nu
(\varepsilon^{\lambda_1},\alpha)^{\nu_1}(\varepsilon^{\lambda_2},\alpha)^{\nu_2}\ldots$ \\
„ — „ 23 v.\,o. „ $\varepsilon^{\lambda_2}$ statt $\varepsilon^2$. \\
„ 551, „ 15 v.\,o. „ $\alpha_k$ „ $\alpha^k$. \\
„ 592, „ 8 v.\,u. „ die nur erf\"ullt ist f\"ur $b=0$, $a=\pm 1$ oder
$a=0$, $b=\pm 1$ oder $a=b=\pm 1$. \\
„ 612, „ 12 v.\,o. „ $1-a$ statt $a-b$. \\
„ 618, „ 1 v.\,u. „ $c_1^h$ statt $c_1$. \\
„ 623, „ 6 v.\,u. „ $a_6$ „ $\alpha_6$.
\\
„ 636, „ 18, 21 v.\,o. lies $-g^{n-1}+1$ statt $g^{n-1}-1$. \\
„ 639, „ 10, 27 v.\,o. „ $n-2$, $n-3$ statt $n-1$, $n-2$. \\
„ 641, „ 1 v.\,o. lies (6) statt (5). \\
„ 646, „ 11 v.\,o. „ rational in $\mathfrak{K}(f_0)$ statt rational. \\
„ — „ 14, 19, 22, 26 v.\,o. lies $\mathfrak{K}(f_0)$ statt $\mathfrak{K}$. \\
„ — „ 17, 18 v.\,o. von \glqq{}Alle Gr\"ossen\grqq{} an zu streichen. \\
„ — „ 1, 2 v.\,u. „ \glqq{}die anderen\grqq{} an zu streichen. \\
„ 653, „ 4 v.\,o. lies $-A_4 r\varrho$ statt $+A_4 r\varrho$. \\
„ — „ 5 v.\,o. „ $+A_2\varrho'$ „ $-A_2\varrho'$.
\end{flushleft}

\begin{center}
\rule{6em}{0.4pt}
\end{center}

\begin{center}
\textbf{\gesperrt{Berichtigungen zum zweiten Bande.}}
\end{center}

\begin{center}
\rule{6em}{0.4pt}
\end{center}

\begin{flushleft}
Seite 15, Zeile 19 v.\,o. lies $N$ statt $P$. \\
„ 16, „ 12, 13 v.\,o. sind $P$ und $Q$ zu vertauschen. \\
„ 21, „ 19 v.\,o. lies (17) statt (13). \\
„ 26, „ 17 v.\,o. „ \S.\,6, 4 statt \S.\,6, 2. \\
„ 569, „ 5 v.\,u. im Exponenten lies $n-\nu$ statt $\mu-\nu$. \\
„ 592, „ 6, 3 v.\,u. lies $\gamma^{P^{f-1}}$ statt $\gamma^{P^f-1}$.
\end{flushleft}

\begin{center}
\rule{6em}{0.4pt}
\end{center}
